\section{Исходный код файла \texttt{Visualizer.h}}
\infoGen{Егоров Александр Владиславович}{swrneko@mai.education}
\footnotesize
\begin{listing}{1}
/****************************************************************
*                     КАФЕДРА № 304 2 КУРС                       *
*---------------------------------------------------------------*
* Project Type  : Qt Widgets Desktop Application                *
* Project Name  : FieldViz                                      *
* File Name     : Visualizer.h                                   *
* Language      : C/C++ (C++17)                                  *
* Programmer(s) : Егоров А.В. (swrneko)                          *
* Modifyed By   : Егоров А.В. (swrneko)                          *
* Edited by     : Neovim                                         *
* OS            : Arch Linux                                     *
* Created       : 01/07/26                                       *
* Last Revision : 01/07/26                                       *
* Comment(s)    : Виджет 3D визуализации полей.                  *
****************************************************************/

#ifndef VISUALIZER_H
#define VISUALIZER_H

#include <QWidget>
#include <QPoint>
#include <vector>
#include "FieldMath.h"

/****************************************************************
* Класс виджета визуализации полей и эквипотенциалей            *
****************************************************************/
class Visualizer : public QWidget {
  Q_OBJECT
public:
  // Конструктор виджета
  explicit Visualizer(QWidget *parent = nullptr);

  // Установка отрисовки силовых линий
  void setShowFieldLines(bool show);

  // Установка отрисовки эквипотенциалей
  void setShowEquipotentials(bool show);

  // Сброс камеры в дефолтное положение
  void resetCamera();

protected:
  // Переопределение события рисования
  void paintEvent(QPaintEvent *event) override;

  // Переопределения мыши для вращения сцены
  void mousePressEvent(QMouseEvent *event) override;
  void mouseMoveEvent(QMouseEvent *event) override;
  void wheelEvent(QWheelEvent *event) override;

private:
  double m_yaw;   // Угол поворота вокруг оси Z (в градусах)
  double m_pitch; // Угол поворота вокруг оси X (в градусах)
  double m_zoom;  // Масштаб отображения
  bool m_showFieldLines; // Флаг показа силовых линий
  bool m_showEquipotentials; // Флаг показа эквипотенциалей

  QPoint m_lastMousePos; // Последняя позиция курсора мыши

  // Функция проекции 3D точки на 2D экран
  QPointF project(const Point3D &p, int width, int height) const;

  // Отрисовка координатных осей
  void drawAxes(QPainter &painter, int w, int h);

  // Отрисовка каркаса единичной сферы (1-й октант)
  void drawSphereWireframe(QPainter &painter, int w, int h);

  // Отрисовка силовых линий векторного поля a
  void drawFieldLines(QPainter &painter, int w, int h);

  // Отрисовка эквипотенциальных линий поля b
  void drawEquipotentials(QPainter &painter, int w, int h);
};

#endif // VISUALIZER_H

\end{listing}
\normalsize
\pagebreak
